\section{The Eight Channels}
\label{sec:two-models}

Four channels carry resource outcomes into the economy, two introduce policy instruments,
and two return economic information to CLEWS. Each scenario uses the combination relevant
to the policy being assessed.

The Philippine illustrations show how the model uses each adjustment and what
changes as a result, before bringing the channels together. They use saved July 2026
experiments, not forecasts. Economic effects are measured against the same model
solved without the channel adjustments, labelled the \emph{control}.\footnote{The
saved unperturbed reform is used throughout to separate channel effects from the
small difference between the original baseline and an unperturbed reform solve.
Return signals retain their reference to the original resource-plan baseline.}
Where a solved response is unavailable, the tables distinguish the model input
from illustrative arithmetic.

The result charts use the same years, scale and colours for GDP, consumption,
capital and labour. This makes small responses visibly small and allows direct
comparison across experiments. Capital denotes private productive capital;
public infrastructure is a separate stock.

\begin{table}[H]
\centering\small
\caption{Economic relationships in OG--CLEWS.}
\label{tab:two-way}
\begin{tabularx}{\linewidth}{@{}>{\bfseries\raggedright\arraybackslash}p{0.28\linewidth} L@{}}
\toprule
Channel & Economic role \\
\midrule
\multicolumn{2}{@{}l}{\emph{Resource outcomes entering OG-Core}} \\
Electricity cost & household purchasing power and industry production costs \\
Public investment & productive infrastructure and its financing \\
Generation capital intensity & the role of capital in electricity production \\
Health & survival, population structure and working capacity \\
\addlinespace
\multicolumn{2}{@{}l}{\emph{Policy instruments}} \\
Carbon price & technology choices, carbon payments and public revenue \\
Private investment incentive & the cost of capital for generation investment \\
\addlinespace
\multicolumn{2}{@{}l}{\emph{Economic feedback to CLEWS}} \\
Planning discount rate & the balance between upfront and future costs \\
Demand & the resource requirements of households and industries \\
\bottomrule
\end{tabularx}
\end{table}
